% GENERATED FILE. Edit canonical lessons, then run npm run generate:books.
% canonical-source-hash: fnv1a64:ccc7aac8b6ad8aab
% canonical-lessons: PA-C184-time-return, PA-C184-hafte, PA-C184-mahine, PA-C184-agli, PA-C184-agle, PA-C184-pichhle, PA-C184-time-review, PA-C184-payoff

\chapter{Next Week, Last Week}
\label{ch:pa-a2-nearby-time}
\hlchaptermodality{184}

\begin{chapteropening}
\textbf{By the end of this chapter:} \emph{I can recognise and use five exact forms for nearby weeks, months, and next or last references.}

Retrieve five exact nearby-time and reference forms through listening, speaking, reading, and bounded copying after a slow review-led ramp.
\end{chapteropening}

\section[(week and month review)]{Start with one week and one month}

\label{lesson:PA-C184-time-return}



\begin{quote}
Read \textbf{\pa{ਹਫ਼ਤਾ}}, then \textbf{\pa{ਮਹੀਨਾ}}.
\end{quote}

\begin{grammarlens}[title={Grammar Lens: keep the two time ideas steady}]
One card names a week. One card names a month. Say each base form twice. The next cards will change only the ending or add one nearby-time word.
\end{grammarlens}

\subsection*{Your turn}
\begin{itemize}
\raggedright
  \item \emph{Hear:} \emph{haftaa}; YOU SAY: \textbf{\pa{ਹਫ਼ਤਾ}}
  \item \emph{Hear:} \emph{mahiinaa}; YOU SAY: \textbf{\pa{ਮਹੀਨਾ}}
  \item \emph{Read it:} \textbf{\pa{ਹਫ਼ਤਾ}, \pa{ਮਹੀਨਾ}}
\end{itemize}

\begin{tcolorbox}[breakable,colback=teal!4,colframe=teal!35!black,title={Before you move on}]
Which card says “a week”? Which says “a month”? (\textbf{\pa{ਹਫ਼ਤਾ}; \pa{ਮਹੀਨਾ}}.)

This review introduces no new written target.
\end{tcolorbox}
\section[\texorpdfstring{\pa{ਹਫ਼ਤੇ} (hafte)}{ਹਫ਼ਤੇ (hafte)}]{\pa{ਹਫ਼ਤੇ} — weeks, or week inside a time phrase}

\label{lesson:PA-C184-hafte}



\begin{quote}
Read the base card \textbf{\pa{ਹਫ਼ਤਾ}}.
\end{quote}

\begin{grammarlens}[title={Grammar Lens: change the final vowel}]
\begin{quote}
\textbf{\pa{ਹਫ਼ਤੇ}} — \emph{hafte} — weeks, or week inside a larger time phrase
\end{quote}

The ending changes from \emph{-aa} to \emph{-e}. In the small phrase \textbf{\pa{ਦੋ} \pa{ਹਫ਼ਤੇ}}, it means “two weeks.”
\end{grammarlens}

\subsection*{Your turn}
\begin{itemize}
\raggedright
  \item \emph{Hear:} \emph{haftaa, hafte}; YOU SAY: \textbf{\pa{ਹਫ਼ਤਾ}, \pa{ਹਫ਼ਤੇ}}
  \item \emph{Read it:} \textbf{\pa{ਦੋ} \pa{ਹਫ਼ਤੇ}}
  \item \emph{Write it:} copy \textbf{\pa{ਹਫ਼ਤੇ}} once and check the final \textbf{\pa{ੇ}}
\end{itemize}

\begin{tcolorbox}[breakable,colback=teal!4,colframe=teal!35!black,title={Before you move on}]
Read the exact form in “two weeks.” (\textbf{\pa{ਹਫ਼ਤੇ}}.)

Source: the canonical Punjabi \textbf{\pa{ਹਫ਼ਤਾ}} entry and the concrete phrase \textbf{\pa{ਦੋ} \pa{ਹਫ਼ਤੇ}}.
\end{tcolorbox}
\section[\texorpdfstring{\pa{ਮਹੀਨੇ} (mahiine)}{ਮਹੀਨੇ (mahiine)}]{\pa{ਮਹੀਨੇ} — month inside a time phrase}

\label{lesson:PA-C184-mahine}



\begin{quote}
Read \textbf{\pa{ਮਹੀਨਾ}}, then yesterday’s changed week form \textbf{\pa{ਹਫ਼ਤੇ}}.
\end{quote}

\begin{grammarlens}[title={Grammar Lens: make the same ending change}]
\begin{quote}
\textbf{\pa{ਮਹੀਨੇ}} — \emph{mahiine} — month inside a larger time phrase
\end{quote}

Again the final \emph{-aa} becomes \emph{-e}. The phrase \textbf{\pa{ਹਰ} \pa{ਮਹੀਨੇ}} means “every month.”
\end{grammarlens}

\subsection*{Your turn}
\begin{itemize}
\raggedright
  \item \emph{Hear:} \emph{mahiinaa, mahiine}; YOU SAY: \textbf{\pa{ਮਹੀਨਾ}, \pa{ਮਹੀਨੇ}}
  \item \emph{Read it:} \textbf{\pa{ਹਰ} \pa{ਮਹੀਨੇ}}
  \item \emph{Write it:} copy \textbf{\pa{ਮਹੀਨੇ}} once and check the final \textbf{\pa{ੇ}}
\end{itemize}

\begin{tcolorbox}[breakable,colback=teal!4,colframe=teal!35!black,title={Before you move on}]
Read the changed week form, then the changed month form. (\textbf{\pa{ਹਫ਼ਤੇ}; \pa{ਮਹੀਨੇ}}.)

Source: the canonical Punjabi \textbf{\pa{ਮਹੀਨਾ}} entry and the concrete phrase \textbf{\pa{ਹਰ} \pa{ਮਹੀਨੇ}}.
\end{tcolorbox}
\section[\texorpdfstring{\pa{ਅਗਲੀ} (aglii)}{ਅਗਲੀ (aglii)}]{\pa{ਅਗਲੀ} — next, with a feminine noun}

\label{lesson:PA-C184-agli}



\begin{quote}
Read \textbf{\pa{ਯਾਤਰਾ}}, then \textbf{\pa{ਮਹੀਨੇ}}.
\end{quote}

\begin{grammarlens}[title={Grammar Lens: one concrete next-item phrase}]
\begin{quote}
\textbf{\pa{ਅਗਲੀ}} — \emph{aglii} — next, with a feminine noun
\end{quote}

Keep it attached to one useful model: \textbf{\pa{ਅਗਲੀ} \pa{ਯਾਤਰਾ}} — “the next journey.”
\end{grammarlens}

\subsection*{Your turn}
\begin{itemize}
\raggedright
  \item \emph{Hear:} \emph{aglii yaatraa}; YOU SAY: \textbf{\pa{ਅਗਲੀ} \pa{ਯਾਤਰਾ}}
  \item \emph{Read it:} \textbf{\pa{ਅਗਲੀ} \pa{ਯਾਤਰਾ}}
  \item \emph{Write it:} copy \textbf{\pa{ਅਗਲੀ}} once, then copy the two-word phrase
\end{itemize}

\begin{tcolorbox}[breakable,colback=teal!4,colframe=teal!35!black,title={Before you move on}]
Read “the next journey.” (\textbf{\pa{ਅਗਲੀ} \pa{ਯਾਤਰਾ}}.)

Source: standard Punjabi agreement in the concrete phrase \textbf{\pa{ਅਗਲੀ} \pa{ਯਾਤਰਾ}}.
\end{tcolorbox}
\section[\texorpdfstring{\pa{ਅਗਲੇ} (agle)}{ਅਗਲੇ (agle)}]{\pa{ਅਗਲੇ} — next, in a masculine or plural phrase}

\label{lesson:PA-C184-agle}



\begin{quote}
Read \textbf{\pa{ਅਗਲੀ}}, then \textbf{\pa{ਮਹੀਨੇ}}.
\end{quote}

\begin{grammarlens}[title={Grammar Lens: change only the agreement ending}]
\begin{quote}
\textbf{\pa{ਅਗਲੇ}} — \emph{agle} — next, in a masculine or plural phrase
\end{quote}

Compare the final sounds \emph{-ii} and \emph{-e}. Keep this model together: \textbf{\pa{ਅਗਲੇ} \pa{ਮਹੀਨੇ}} — “next month.”
\end{grammarlens}

\subsection*{Your turn}
\begin{itemize}
\raggedright
  \item \emph{Hear:} \emph{aglii, agle}; YOU SAY: \textbf{\pa{ਅਗਲੀ}, \pa{ਅਗਲੇ}}
  \item \emph{Read it:} \textbf{\pa{ਅਗਲੇ} \pa{ਮਹੀਨੇ}}
  \item \emph{Write it:} copy \textbf{\pa{ਅਗਲੇ} \pa{ਮਹੀਨੇ}} once and check both final \textbf{\pa{ੇ}} signs
\end{itemize}

\begin{tcolorbox}[breakable,colback=teal!4,colframe=teal!35!black,title={Before you move on}]
Read “the next journey,” then “next month.” (\textbf{\pa{ਅਗਲੀ} \pa{ਯਾਤਰਾ}; \pa{ਅਗਲੇ} \pa{ਮਹੀਨੇ}}.)

Source: standard Punjabi agreement in the concrete phrase \textbf{\pa{ਅਗਲੇ} \pa{ਮਹੀਨੇ}}.
\end{tcolorbox}
\section[\texorpdfstring{\pa{ਪਿਛਲੇ} (pichhle)}{ਪਿਛਲੇ (pichhle)}]{\pa{ਪਿਛਲੇ} — last or previous in a time phrase}

\label{lesson:PA-C184-pichhle}



\begin{quote}
Read \textbf{\pa{ਅਗਲੇ}}, then \textbf{\pa{ਹਫ਼ਤੇ}}.
\end{quote}

\begin{grammarlens}[title={Grammar Lens: point backward instead of forward}]
\begin{quote}
\textbf{\pa{ਪਿਛਲੇ}} — \emph{pichhle} — last or previous, in a masculine time phrase
\end{quote}

Keep one compact contrast: \textbf{\pa{ਅਗਲੇ} \pa{ਮਹੀਨੇ}} points forward; \textbf{\pa{ਪਿਛਲੇ} \pa{ਹਫ਼ਤੇ}} points backward.
\end{grammarlens}

\subsection*{Your turn}
\begin{itemize}
\raggedright
  \item \emph{Hear:} \emph{agle, pichhle}; YOU SAY: \textbf{\pa{ਅਗਲੇ}, \pa{ਪਿਛਲੇ}}
  \item \emph{Read it:} \textbf{\pa{ਪਿਛਲੇ} \pa{ਹਫ਼ਤੇ}}
  \item \emph{Write it:} copy \textbf{\pa{ਪਿਛਲੇ} \pa{ਹਫ਼ਤੇ}} once
\end{itemize}

\begin{tcolorbox}[breakable,colback=teal!4,colframe=teal!35!black,title={Before you move on}]
Read the backward phrase, then the forward phrase. (\textbf{\pa{ਪਿਛਲੇ} \pa{ਹਫ਼ਤੇ}; \pa{ਅਗਲੇ} \pa{ਮਹੀਨੇ}}.)

Source: standard Punjabi time reference in the concrete phrase \textbf{\pa{ਪਿਛਲੇ} \pa{ਹਫ਼ਤੇ}}.
\end{tcolorbox}
\section[\texorpdfstring{\pa{ਹਫ਼ਤੇ}, \pa{ਮਹੀਨੇ}, \pa{ਅਗਲੀ} … (hafte, mahiine …)}{ਹਫ਼ਤੇ, ਮਹੀਨੇ, ਅਗਲੀ … (hafte, mahiine …)}]{Five forms in five small phrases}

\label{lesson:PA-C184-time-review}



\begin{quote}
Read \textbf{\pa{ਦੋ} \pa{ਹਫ਼ਤੇ}}, then \textbf{\pa{ਹਰ} \pa{ਮਹੀਨੇ}}.
\end{quote}

\subsection*{Your turn: keep each form in its model}
Read slowly:

\begin{itemize}
\raggedright
  \item \textbf{\pa{ਦੋ} \pa{ਹਫ਼ਤੇ}} — two weeks
  \item \textbf{\pa{ਹਰ} \pa{ਮਹੀਨੇ}} — every month
  \item \textbf{\pa{ਅਗਲੀ} \pa{ਯਾਤਰਾ}} — the next journey
  \item \textbf{\pa{ਅਗਲੇ} \pa{ਮਹੀਨੇ}} — next month
  \item \textbf{\pa{ਪਿਛਲੇ} \pa{ਹਫ਼ਤੇ}} — last week
\end{itemize}

\emph{Hear:} the five phrases in mixed order; YOU SAY: each matching time meaning

\begin{tcolorbox}[breakable,colback=teal!4,colframe=teal!35!black,title={Before you move on}]
\emph{Write it:} cover the models and write \textbf{\pa{ਹਫ਼ਤੇ}, \pa{ਮਹੀਨੇ}, \pa{ਅਗਲੀ}, \pa{ਅਗਲੇ}, \pa{ਪਿਛਲੇ}} once

Check the two time nouns first, then the three nearby-time forms.

This review adds no new form and makes no mock-readiness claim.
\end{tcolorbox}
\section[\texorpdfstring{\pa{ਪਿਛਲੇ} \pa{ਹਫ਼ਤੇ} · \pa{ਅਗਲੀ} … (pichhle hafte …)}{ਪਿਛਲੇ ਹਫ਼ਤੇ · ਅਗਲੀ … (pichhle hafte …)}]{A gentle nearby-time payoff}

\label{lesson:PA-C184-payoff}



\begin{quote}
Without looking, say the changed forms for week and month.
\end{quote}

\subsection*{Your turn: four skills, bounded to five forms}
\begin{itemize}
\raggedright
  \item \emph{Hear:} \emph{pichhle hafte, aglii yaatraa, agle mahiine}; YOU SAY: \textbf{\pa{ਪਿਛਲੇ} \pa{ਹਫ਼ਤੇ}, \pa{ਅਗਲੀ} \pa{ਯਾਤਰਾ}, \pa{ਅਗਲੇ} \pa{ਮਹੀਨੇ}}
  \item \emph{Read it:} \textbf{\pa{ਦੋ} \pa{ਹਫ਼ਤੇ।} \pa{ਹਰ} \pa{ਮਹੀਨੇ।}}
  \item \emph{Write it:} cover the model and write \textbf{\pa{ਹਫ਼ਤੇ}, \pa{ਮਹੀਨੇ}, \pa{ਅਗਲੀ}, \pa{ਅਗਲੇ}, \pa{ਪਿਛਲੇ}}
\end{itemize}

Now point to the form that agrees with \textbf{\pa{ਯਾਤਰਾ}}. Then point to the two forms that end in \textbf{\pa{ਲੇ}}.

\begin{tcolorbox}[breakable,colback=teal!4,colframe=teal!35!black,title={Before you move on}]
Retrieve all five exact forms once more in any order.

This is the seventh prerequisite tranche exposed by Punjabi A2 mock drafting. It does not unblock Punjabi A2 mock 1 by itself. It is not an A2 mock, not A2 readiness or attainment, and not external or human validation.

Source: the canonical Punjabi week, month, journey, and practical-schedule entries.
\end{tcolorbox}
